\begin{tikzpicture}\begin{axis}[otLivreGraphe,width=95mm,height=47mm,
xmin=-0.05,xmax=1.1,ymin=-1.1,ymax=1.1,xtick={0,.5,1},ytick={-1,0,1},scaled ticks=false,
tick label style={/pgf/number format/use comma},clip=true]

\addplot[otNavy,thick,domain=0:1,samples=90]{x^3+x-1};
\draw[otGreenDark,dashed](axis cs:.682327804,0)--(axis cs:.682327804,-1);
\node[above left,font=\small]at(axis cs:.682327804,0){$\alpha$};

\end{axis}\end{tikzpicture}
\par\smallskip\legende{Le graphique situe la racine ; les signes aux bornes et la largeur de l’intervalle donnent l’encadrement garanti.}
